\documentclass[10pt,a4paper]{amsart}
\usepackage[margin=19mm]{geometry}
\usepackage{amsmath,amssymb,mathtools}
\usepackage[scr=rsfs,cal=euler]{mathalfa}
\usepackage{xfrac}
\usepackage[unicode,hidelinks,bookmarks=false]{hyperref}
\newcommand{\ie}{i.e.\ }
\newcommand{\cf}{cf.\ }
\newcommand{\resp}{respectively\ }
\newcommand{\Subsection}{\subsection}
\providecommand{\nobreakdash}{\nobreak\hbox{-}}
\setlength{\emergencystretch}{2em}
\multlinegap0pt
\usepackage{tikz}
\usetikzlibrary{cd}
\newtheorem{prop}{Proposition}
\newtheorem{defi}{Definition}
\usepackage{cleveref}
\crefname{prop}{Proposition}{Propositions}
\crefname{defi}{Definition}{Definitions}
\newcommand{\dbl}[1]{\ifstrequal{#1}{1}{\mathds{1}}{\mathbb{#1}}}
\DeclareMathOperator{\id}{id}
\def\p#1{\mathrel{\ooalign{\hfil$\mapstochar\mkern 5mu$\hfil\cr$#1$}}}
\newcommand{\proto}{\p\rightarrow}
\title{Quantification in Double-Categorical Database Schemas}
\author{Michael Lambert}
\date{Source: arXiv:2608.00913v2 (CC BY 4.0)}
\begin{document}
\maketitle

\noindent\emph{Source and licence.} Quantification in Double-Categorical Database Schemas. Version arXiv:2608.00913v2 (CC BY 4.0), distributed by arXiv under the Creative Commons Attribution 4.0 International licence, http://creativecommons.org/licenses/by/4.0/, as recorded in the arXiv licence metadata for this exact version and shown on the abstract page https://arxiv.org/abs/2608.00913v2. Full licence notice: \texttt{licenses/arxiv-2608.00913-CC-BY-4.0.txt}.\par\smallskip

\noindent\textit{Editorial scope.} Counted unit: the article's prop prop:modal-properties (source0.tex lines 1624--1637). The article's complete original proof of it is delivered with it (source0.tex lines 1638--1694, one \texttt{proof} environment of 57 lines). No mathematics is written, repaired or re-derived: the statement and the proof are the article's own lines, in the article's own notation, and no retained line is substituted or re-flowed.  Retained uncounted as the source-local context the proof needs: lines 869--904; lines 910--912; lines 1592--1594; lines 1602--1610.  Nothing was omitted from any retained span, so the package carries no omission record.  No retained line contains a citation, so the wrapper carries no bibliography and no bibliography entry.  The retained lines contain the article's own diagram code, which the wrapper typesets with the standard tikz packages; no image file is delivered and the package declares no asset dependency.  Editorial formatting only, in addition: an AMS \texttt{amsart} preamble that restates the article's own declarations and macros used by the retained lines, namely the environments \texttt{prop}, \texttt{defi} on separate counters, together with the standard packages those lines need.  The retained environments print in this document as Proposition 1, Definition 1 in order of appearance, whereas the article prints the same items with the numbering of its own full document.  The complete native source of the article as distributed in the arXiv e-print for this exact version is bundled unchanged as \texttt{sources/206579/source0.tex}.
    \begin{equation} \label{equation:globular-cell}
      % https://q.uiver.app/#q=WzAsOCxbMCwxLCJ4Il0sWzEsMSwicyJdLFsyLDEsInMiXSxbMiwyLCIxIl0sWzAsMiwieCJdLFsxLDIsInkiXSxbMCwwLCJ4Il0sWzIsMCwicyJdLFsxLDIsIlxcaWQiLDAseyJzdHlsZSI6eyJib2R5Ijp7Im5hbWUiOiJiYXJyZWQifX19XSxbMiwzLCIhIl0sWzAsNCwiMSIsMix7ImxldmVsIjoyLCJzdHlsZSI6eyJoZWFkIjp7Im5hbWUiOiJub25lIn19fV0sWzQsNSwiciIsMix7InN0eWxlIjp7ImJvZHkiOnsibmFtZSI6ImJhcnJlZCJ9fX1dLFs1LDMsIlxcaGF0XFxwaGkiLDIseyJzdHlsZSI6eyJib2R5Ijp7Im5hbWUiOiJiYXJyZWQifX19XSxbMSw1LCJcXHBoaSJdLFs2LDAsIiIsMCx7ImxldmVsIjoyLCJzdHlsZSI6eyJoZWFkIjp7Im5hbWUiOiJub25lIn19fV0sWzYsNywiclxcb2RvdCBcXHBoaV4qIiwwLHsic3R5bGUiOnsiYm9keSI6eyJuYW1lIjoiYmFycmVkIn19fV0sWzcsMl0sWzAsMSwiclxcb2RvdCBcXHBoaV4qIiwwLHsic3R5bGUiOnsiYm9keSI6eyJuYW1lIjoiYmFycmVkIn19fV0sWzgsMTIsIlxcdGV4dHtleHR9IiwxLHsic2hvcnRlbiI6eyJzb3VyY2UiOjIwLCJ0YXJnZXQiOjIwfSwic3R5bGUiOnsiYm9keSI6eyJuYW1lIjoibm9uZSJ9LCJoZWFkIjp7Im5hbWUiOiJub25lIn19fV0sWzE3LDExLCJcXHRleHR7cmVzfSIsMSx7InNob3J0ZW4iOnsic291cmNlIjoyMCwidGFyZ2V0IjoyMH0sInN0eWxlIjp7ImJvZHkiOnsibmFtZSI6Im5vbmUifSwiaGVhZCI6eyJuYW1lIjoibm9uZSJ9fX1dLFsxNSwxLCJcXGNvbmciLDEseyJzaG9ydGVuIjp7InNvdXJjZSI6MjB9LCJzdHlsZSI6eyJib2R5Ijp7Im5hbWUiOiJub25lIn0sImhlYWQiOnsibmFtZSI6Im5vbmUifX19XV0=
      \begin{tikzcd}
        x && s \\
        x & s & s \\
        x & y & 1
        \arrow[""{name=0, anchor=center, inner sep=0}, "{r\odot \phi^*}"{inner sep=.8ex}, "\shortmid"{marking}, from=1-1, to=1-3]
        \arrow[equals, from=1-1, to=2-1]
        \arrow[from=1-3, to=2-3]
        \arrow[""{name=1, anchor=center, inner sep=0}, "{r\odot \phi^*}"{inner sep=.8ex}, "\shortmid"{marking}, from=2-1, to=2-2]
        \arrow["1"', equals, from=2-1, to=3-1]
        \arrow[""{name=2, anchor=center, inner sep=0}, "\id"{inner sep=.8ex}, "\shortmid"{marking}, from=2-2, to=2-3]
        \arrow["\phi", from=2-2, to=3-2]
        \arrow["{!}", from=2-3, to=3-3]
        \arrow[""{name=3, anchor=center, inner sep=0}, "r"'{inner sep=.8ex}, "\shortmid"{marking}, from=3-1, to=3-2]
        \arrow[""{name=4, anchor=center, inner sep=0}, "{\hat\phi}"'{inner sep=.8ex}, "\shortmid"{marking}, from=3-2, to=3-3]
        \arrow["\cong"{description}, draw=none, from=0, to=2-2]
        \arrow["{\text{res}}"{description}, draw=none, from=1, to=3]
        \arrow["{\text{ext}}"{description}, draw=none, from=2, to=4]
      \end{tikzcd} \qquad = \qquad 
      % https://q.uiver.app/#q=WzAsNixbMCwwLCJ4Il0sWzEsMCwicyJdLFsxLDEsIjEiXSxbMCwxLCJ4Il0sWzAsMiwieCJdLFsxLDIsIjEiXSxbMCwxLCJyXFxvZG90IFxccGhpXioiLDAseyJzdHlsZSI6eyJib2R5Ijp7Im5hbWUiOiJiYXJyZWQifX19XSxbMSwyLCIhIl0sWzAsMywiIiwyLHsibGV2ZWwiOjIsInN0eWxlIjp7ImhlYWQiOnsibmFtZSI6Im5vbmUifX19XSxbMywyLCJcXGRpYW1vbmRzdWl0X3IgXFxwaGkiLDIseyJzdHlsZSI6eyJib2R5Ijp7Im5hbWUiOiJiYXJyZWQifX19XSxbMyw0LCIiLDIseyJsZXZlbCI6Miwic3R5bGUiOnsiaGVhZCI6eyJuYW1lIjoibm9uZSJ9fX1dLFs0LDUsInJcXG9kb3QgXFxoYXRcXHBoaSIsMix7InN0eWxlIjp7ImJvZHkiOnsibmFtZSI6ImJhcnJlZCJ9fX1dLFsyLDUsIiIsMix7ImxldmVsIjoyLCJzdHlsZSI6eyJoZWFkIjp7Im5hbWUiOiJub25lIn19fV0sWzksMTEsIlxcZXhpc3RzXFwsISIsMSx7ImxhYmVsX3Bvc2l0aW9uIjo2MCwic2hvcnRlbiI6eyJzb3VyY2UiOjIwLCJ0YXJnZXQiOjIwfSwic3R5bGUiOnsiYm9keSI6eyJuYW1lIjoibm9uZSJ9LCJoZWFkIjp7Im5hbWUiOiJub25lIn19fV0sWzYsOSwiXFx0ZXh0e2V4dH0iLDEseyJzaG9ydGVuIjp7InNvdXJjZSI6MjAsInRhcmdldCI6MjB9LCJzdHlsZSI6eyJib2R5Ijp7Im5hbWUiOiJub25lIn0sImhlYWQiOnsibmFtZSI6Im5vbmUifX19XV0=
      \begin{tikzcd}
        x & s \\
        x & 1 \\
        x & 1
        \arrow[""{name=0, anchor=center, inner sep=0}, "{r\odot \phi^*}"{inner sep=.8ex}, "\shortmid"{marking}, from=1-1, to=1-2]
        \arrow[equals, from=1-1, to=2-1]
        \arrow["{!}", from=1-2, to=2-2]
        \arrow[""{name=1, anchor=center, inner sep=0}, "{\diamondsuit_r \phi}"'{inner sep=.8ex}, "\shortmid"{marking}, from=2-1, to=2-2]
        \arrow[equals, from=2-1, to=3-1]
        \arrow[equals, from=2-2, to=3-2]
        \arrow[""{name=2, anchor=center, inner sep=0}, "{r\odot \hat\phi}"'{inner sep=.8ex}, "\shortmid"{marking}, from=3-1, to=3-2]
        \arrow["{\text{ext}}"{description}, draw=none, from=0, to=1]
        \arrow["{\exists\,!}"{description, pos=0.6}, draw=none, from=1, to=2]
      \end{tikzcd}.
    \end{equation}

\begin{prop} \label{prop:possibility-rewrite-rule-1}
  The globular cell $\diamondsuit_r \phi \Rightarrow r\odot \hat\phi$ induced in \cref{equation:globular-cell} is an isomorphism in any equipment. Accordingly, the two potential definitions of possibility agree.
\end{prop}

  \begin{equation} \label{eqn:possibility-operator}
    \diamondsuit_r:=r\odot - \colon \dbl D(y,1)\to\dbl D(x,1) \qquad\qquad p\colon y\proto 1 \quad \mapsto \quad r\odot p
  \end{equation}

\begin{defi}[Adjoint Modality] \label{def:adjoint-modality}
  The operator $\diamondsuit_r$ of \cref{eqn:possibility-operator}
  is the left adjoint of an \textbf{adjoint modality} if it has a 
  right adjoint 
    \begin{equation*}
      \Box\colon\dbl D(x,1)\to\dbl D(y,1) \qquad \diamondsuit_r\dashv\Box
    \end{equation*}
  called \emph{(reverse) necessity}. The pair is referred to as an adjoint modality where $\Box$ is the right adjoint of the modality.
\end{defi}

\begin{prop} \label{prop:modal-properties}
  Let $r\colon x\proto y$ induce an adjoint modality 
  (\cref{def:adjoint-modality}). If $r$ is a preorder, then the possibility and 
  necessity operators associated to $r$ satisfy the monad and comonad axioms 
    \begin{enumerate}
      \item $\id\leq\diamondsuit$ and $\diamondsuit\diamondsuit\leq\diamondsuit$
      \item $\Box\leq \id$ and $\Box\Box\leq\Box$
    \end{enumerate}
  as well as the further (B) axioms, namely, 
    \begin{enumerate}
      \item $\id\leq \Box\diamondsuit$
      \item $\diamondsuit\Box\leq \id$.
    \end{enumerate}
\end{prop}

\begin{proof}
  Since $\diamondsuit\dashv\Box$, it suffices to show that $\diamondsuit$ 
  satisfies the monad axioms as $\Box$ will automatically be a comonad. Likewise 
  the (B) axioms are just a restatement of the adjunction assumption. 
  Reflexivity of $r$ means that there is a canonical cell 
    \begin{equation*}
      % https://q.uiver.app/#q=WzAsNCxbMCwwLCJ4Il0sWzEsMCwieCJdLFsxLDEsIngiXSxbMCwxLCJ4Il0sWzAsMSwiXFxpZF94IiwwLHsic3R5bGUiOnsiYm9keSI6eyJuYW1lIjoiYmFycmVkIn19fV0sWzEsMiwiIiwwLHsibGV2ZWwiOjIsInN0eWxlIjp7ImhlYWQiOnsibmFtZSI6Im5vbmUifX19XSxbMCwzLCIiLDIseyJsZXZlbCI6Miwic3R5bGUiOnsiaGVhZCI6eyJuYW1lIjoibm9uZSJ9fX1dLFszLDIsInIiLDIseyJzdHlsZSI6eyJib2R5Ijp7Im5hbWUiOiJiYXJyZWQifX19XSxbNCw3LCJcXGRlbHRhX3giLDEseyJzaG9ydGVuIjp7InNvdXJjZSI6MjAsInRhcmdldCI6MjB9LCJzdHlsZSI6eyJib2R5Ijp7Im5hbWUiOiJub25lIn0sImhlYWQiOnsibmFtZSI6Im5vbmUifX19XV0=
      \begin{tikzcd}
        x & x \\
        x & x
        \arrow[""{name=0, anchor=center, inner sep=0}, "{\id_x}"{inner sep=.8ex}, "\shortmid"{marking}, from=1-1, to=1-2]
        \arrow[equals, from=1-1, to=2-1]
        \arrow[equals, from=1-2, to=2-2]
        \arrow[""{name=1, anchor=center, inner sep=0}, "r"'{inner sep=.8ex}, "\shortmid"{marking}, from=2-1, to=2-2]
        \arrow["{\delta_x}"{description}, draw=none, from=0, to=1]
      \end{tikzcd}
    \end{equation*}
  expressing the fact that the diagonal factors through $r$. But then composing 
  we have 
    \begin{equation*}
      % https://q.uiver.app/#q=WzAsNixbMCwwLCIxIl0sWzEsMCwieCJdLFsxLDEsIngiXSxbMCwxLCIxIl0sWzIsMCwieCJdLFsyLDEsIngiXSxbMCwxLCJwIiwwLHsic3R5bGUiOnsiYm9keSI6eyJuYW1lIjoiYmFycmVkIn19fV0sWzEsMiwiIiwwLHsibGV2ZWwiOjIsInN0eWxlIjp7ImhlYWQiOnsibmFtZSI6Im5vbmUifX19XSxbMCwzLCIiLDIseyJsZXZlbCI6Miwic3R5bGUiOnsiaGVhZCI6eyJuYW1lIjoibm9uZSJ9fX1dLFszLDIsInAiLDIseyJzdHlsZSI6eyJib2R5Ijp7Im5hbWUiOiJiYXJyZWQifX19XSxbMSw0LCJcXGlkIiwwLHsic3R5bGUiOnsiYm9keSI6eyJuYW1lIjoiYmFycmVkIn19fV0sWzIsNSwiUiIsMix7InN0eWxlIjp7ImJvZHkiOnsibmFtZSI6ImJhcnJlZCJ9fX1dLFs0LDUsIiIsMSx7ImxldmVsIjoyLCJzdHlsZSI6eyJoZWFkIjp7Im5hbWUiOiJub25lIn19fV0sWzYsOSwiMV9wIiwxLHsic2hvcnRlbiI6eyJzb3VyY2UiOjIwLCJ0YXJnZXQiOjIwfSwic3R5bGUiOnsiYm9keSI6eyJuYW1lIjoibm9uZSJ9LCJoZWFkIjp7Im5hbWUiOiJub25lIn19fV0sWzEwLDExLCJcXGxlcSIsMSx7InNob3J0ZW4iOnsic291cmNlIjoyMCwidGFyZ2V0IjoyMH0sInN0eWxlIjp7ImJvZHkiOnsibmFtZSI6Im5vbmUifSwiaGVhZCI6eyJuYW1lIjoibm9uZSJ9fX1dXQ==
      \begin{tikzcd}
        1 & x & x \\
        1 & x & x
        \arrow[""{name=0, anchor=center, inner sep=0}, "p"{inner sep=.8ex}, "\shortmid"{marking}, from=1-1, to=1-2]
        \arrow[equals, from=1-1, to=2-1]
        \arrow[""{name=1, anchor=center, inner sep=0}, "\id"{inner sep=.8ex}, "\shortmid"{marking}, from=1-2, to=1-3]
        \arrow[equals, from=1-2, to=2-2]
        \arrow[equals, from=1-3, to=2-3]
        \arrow[""{name=2, anchor=center, inner sep=0}, "p"'{inner sep=.8ex}, "\shortmid"{marking}, from=2-1, to=2-2]
        \arrow[""{name=3, anchor=center, inner sep=0}, "r"'{inner sep=.8ex}, "\shortmid"{marking}, from=2-2, to=2-3]
        \arrow["{1_p}"{description}, draw=none, from=0, to=2]
        \arrow["\leq"{description}, draw=none, from=1, to=3]
      \end{tikzcd}
    \end{equation*}
  showing that $p\leq \diamondsuit p$ by loose composition and 
  \cref{prop:possibility-rewrite-rule-1}. Likewise, using reflexivity and the 
  same proposition, we have 
    \begin{equation*}
      % https://q.uiver.app/#q=WzAsNyxbMCwwLCIxIl0sWzEsMCwieCJdLFsxLDEsIngiXSxbMCwxLCIxIl0sWzIsMCwieCJdLFszLDEsIngiXSxbMywwLCJ4Il0sWzAsMSwicCIsMCx7InN0eWxlIjp7ImJvZHkiOnsibmFtZSI6ImJhcnJlZCJ9fX1dLFsxLDIsIiIsMCx7ImxldmVsIjoyLCJzdHlsZSI6eyJoZWFkIjp7Im5hbWUiOiJub25lIn19fV0sWzAsMywiIiwyLHsibGV2ZWwiOjIsInN0eWxlIjp7ImhlYWQiOnsibmFtZSI6Im5vbmUifX19XSxbMywyLCJwIiwyLHsic3R5bGUiOnsiYm9keSI6eyJuYW1lIjoiYmFycmVkIn19fV0sWzEsNCwiciIsMCx7InN0eWxlIjp7ImJvZHkiOnsibmFtZSI6ImJhcnJlZCJ9fX1dLFsyLDUsInIiLDIseyJzdHlsZSI6eyJib2R5Ijp7Im5hbWUiOiJiYXJyZWQifX19XSxbNCw2LCJyIiwwLHsic3R5bGUiOnsiYm9keSI6eyJuYW1lIjoiYmFycmVkIn19fV0sWzYsNSwiIiwwLHsibGV2ZWwiOjIsInN0eWxlIjp7ImhlYWQiOnsibmFtZSI6Im5vbmUifX19XSxbNywxMCwiMV9wIiwxLHsic2hvcnRlbiI6eyJzb3VyY2UiOjIwLCJ0YXJnZXQiOjIwfSwic3R5bGUiOnsiYm9keSI6eyJuYW1lIjoibm9uZSJ9LCJoZWFkIjp7Im5hbWUiOiJub25lIn19fV0sWzQsMTIsIlxcbGVxIiwxLHsic2hvcnRlbiI6eyJzb3VyY2UiOjIwLCJ0YXJnZXQiOjIwfSwic3R5bGUiOnsiYm9keSI6eyJuYW1lIjoibm9uZSJ9LCJoZWFkIjp7Im5hbWUiOiJub25lIn19fV1d
      \begin{tikzcd}
        1 & x & x & x \\
        1 & x && x
        \arrow[""{name=0, anchor=center, inner sep=0}, "p"{inner sep=.8ex}, "\shortmid"{marking}, from=1-1, to=1-2]
        \arrow[equals, from=1-1, to=2-1]
        \arrow["r"{inner sep=.8ex}, "\shortmid"{marking}, from=1-2, to=1-3]
        \arrow[equals, from=1-2, to=2-2]
        \arrow["r"{inner sep=.8ex}, "\shortmid"{marking}, from=1-3, to=1-4]
        \arrow[equals, from=1-4, to=2-4]
        \arrow[""{name=1, anchor=center, inner sep=0}, "p"'{inner sep=.8ex}, "\shortmid"{marking}, from=2-1, to=2-2]
        \arrow[""{name=2, anchor=center, inner sep=0}, "r"'{inner sep=.8ex}, "\shortmid"{marking}, from=2-2, to=2-4]
        \arrow["{1_p}"{description}, draw=none, from=0, to=1]
        \arrow["\leq"{description}, draw=none, from=1-3, to=2]
      \end{tikzcd}
    \end{equation*}
  so that $\diamondsuit \diamondsuit p \leq \diamondsuit p$ holds, as required.
\end{proof}

\end{document}
